\documentclass[11pt,a4paper]{ltjsarticle}
\usepackage{amsmath,amssymb}
\usepackage{graphicx}
\usepackage{booktabs}
\usepackage{array}
\usepackage{luatexja-fontspec}
\setmainjfont{HaranoAjiMincho-Regular}
\setsansjfont{HaranoAjiGothic-Medium}

% Recommended PDF filename:
% 2611193_HeizoNakai_visual_media_processing1_dai1kai.pdf
% Example build command:
% latexmk -pdfdvi -jobname=2611193_HeizoNakai_visual_media_processing1_dai1kai main.tex

\begin{document}

\begin{center}
{\LARGE \textbf{視覚メディア処理I 第1回レポート}}
\end{center}
\vspace{1em}

\begin{flushright}
提出日: \today \\
学籍番号: 2611193 \\
氏名: 中井 平蔵
\end{flushright}

\section*{課題}
\begin{enumerate}
  \item スライド24ページの回転表現（オイラー角，回転行列，Rodrigues 表現，クォータニオン）の相互変換方法を説明する。
  \item 次のカメラに対して透視投影行列を求める。
  \begin{itemize}
    \item 焦点距離: $f=50\ \mathrm{mm}$
    \item センササイズ: $33.6\ \mathrm{mm}\times 22.4\ \mathrm{mm}$
    \item 画素数: $6720\times 4480\ \mathrm{pixel}$
    \item 視線軸はセンサ中心を通る（主点は画像中心）
  \end{itemize}
\end{enumerate}

\section*{解答}
\subsection*{1. 回転表現の相互変換}
\setlength{\parindent}{1em}
3次元回転は同じ幾何学的意味を，複数の表現で書き換えられる。実装では
「どの表現を中間表現にするか」を固定すると変換を整理しやすい。ここでは
回転行列 $R$ を中心として説明する。

\subsubsection*{(0) 各表現の定義}
\begin{itemize}
  \item オイラー角: 3つの角 $(\phi,\theta,\psi)$ で回転を表す。回転順序（XYZ、ZYXなど）を必ず指定する。
  \item 回転行列: 直交行列 $R\in SO(3)$（$R^\top R=I,\ \det R=1$）。
  \item Rodrigues 表現（軸角）: 単位軸 $\mathbf{u}$ と角度 $\theta$ の組 $(\mathbf{u},\theta)$。
  \item クォータニオン: 単位クォータニオン $q=(w,x,y,z),\ \|q\|=1$。
\end{itemize}
以降では，ベクトルは列ベクトルとする。

\subsubsection*{(a) 3×3回転行列 $\leftrightarrow$ オイラー角}
\hspace*{1em}\textbf{(a-1) オイラー角 $\rightarrow$ 回転行列}

回転順序を ZYX（yaw-pitch-roll）と固定し，角度を
$\psi$（$z$軸），$\theta$（$y$軸），$\phi$（$x$軸）とする。
\[
R_x(\phi)=
\begin{bmatrix}
1&0&0\\
0&\cos\phi&-\sin\phi\\
0&\sin\phi&\cos\phi
\end{bmatrix},\ 
R_y(\theta)=
\begin{bmatrix}
\cos\theta&0&\sin\theta\\
0&1&0\\
-\sin\theta&0&\cos\theta
\end{bmatrix},\ 
R_z(\psi)=
\begin{bmatrix}
\cos\psi&-\sin\psi&0\\
\sin\psi&\cos\psi&0\\
0&0&1
\end{bmatrix}
\]
とすると，
\[
R = R_z(\psi)R_y(\theta)R_x(\phi)
\]
で回転行列を得る。

\hspace*{1em}\textbf{(a-2) 回転行列 $\rightarrow$ オイラー角}

逆変換は $R=[r_{ij}]$ から
\[
\theta = \arcsin(-r_{31}),\quad
\phi = \mathrm{atan2}(r_{32},r_{33}),\quad
\psi = \mathrm{atan2}(r_{21},r_{11})
\]
で求める（この式は ZYX 順専用）。注意として，$\cos\theta\approx 0$
（$\theta\approx\pm\frac{\pi}{2}$）ではジンバルロックが起こり，角度 $\phi$ と $\psi$
を一意に分離できないため，別の式や別表現への切替が必要である。

\subsubsection*{(b) 回転行列 $\leftrightarrow$ Rodrigues 表現（軸角）}
回転軸の単位ベクトルを $\mathbf{u}=(u_x,u_y,u_z)^\top$，回転角を $\theta$ とすると，
\[
R = I + \sin\theta\,[\mathbf{u}]_\times + (1-\cos\theta)[\mathbf{u}]_\times^2
\]
（Rodrigues の回転公式）である。ここで
\[
[\mathbf{u}]_\times=
\begin{bmatrix}
0 & -u_z & u_y\\
u_z & 0 & -u_x\\
-u_y & u_x & 0
\end{bmatrix}.
\]
逆に $R$ からは
\[
\theta=\cos^{-1}\!\left(\frac{\mathrm{tr}(R)-1}{2}\right),\quad
\mathbf{u}=\frac{1}{2\sin\theta}
\begin{bmatrix}
r_{32}-r_{23}\\
r_{13}-r_{31}\\
r_{21}-r_{12}
\end{bmatrix}
\]
で軸角を復元できる（$\theta\approx 0$ では近似が必要）。
実装ではまず $\theta$ を計算し，$\theta$ が十分小さいときは
$\mathbf{u}$ を前回値で保持するか，
$\mathbf{u}\theta$（回転ベクトル）で直接扱うと安定する。

\subsubsection*{(c) 回転行列 $\leftrightarrow$ クォータニオン}
単位クォータニオンを
\[
q=(w,x,y,z),\quad \|q\|=1
\]
とすると，対応する回転行列は
\[
R=
\begin{bmatrix}
1-2(y^2+z^2) & 2(xy-wz) & 2(xz+wy)\\
2(xy+wz) & 1-2(x^2+z^2) & 2(yz-wx)\\
2(xz-wy) & 2(yz+wx) & 1-2(x^2+y^2)
\end{bmatrix}.
\]
逆変換は $\mathrm{tr}(R)$ を用いて
\[
w=\frac{1}{2}\sqrt{1+r_{11}+r_{22}+r_{33}}
\]
を起点に，他成分を差分項から求める。例えば $w\neq 0$ なら
\[
x=\frac{r_{32}-r_{23}}{4w},\ 
y=\frac{r_{13}-r_{31}}{4w},\ 
z=\frac{r_{21}-r_{12}}{4w}.
\]
ただし数値安定のため，実装では「$r_{11}$ と $r_{22}$ と $r_{33}$ のうち最大の対角成分」
を基準に分岐して計算する方法が広く使われる。最後に
$q\leftarrow q/\|q\|$ と正規化する。

\subsubsection*{(d) Rodrigues 表現 $\leftrightarrow$ クォータニオン}
軸角 $(\mathbf{u},\theta)$ から
\[
q=
\left(\cos\frac{\theta}{2},\ 
u_x\sin\frac{\theta}{2},\ 
u_y\sin\frac{\theta}{2},\ 
u_z\sin\frac{\theta}{2}\right)
\]
でクォータニオンに変換できる。逆に
\[
\theta=2\arccos(w),\quad
\mathbf{u}=\frac{1}{\sin(\theta/2)}(x,y,z)^\top
\]
で軸角に戻せる（$\sin(\theta/2)\approx 0$ では軸は任意に近い）。

\subsubsection*{(e) 実用上の変換フロー}
直接式を多数覚えるより，次のように $R$ を中間に置くと整理しやすい。
\begin{itemize}
  \item オイラー角 $\rightarrow R \rightarrow q$（姿勢推定・補間でよく使う）
  \item $q \rightarrow R \rightarrow$ オイラー角（表示用の角度に戻す）
  \item Rodrigues $\rightarrow R \rightarrow q$（最適化変数から姿勢へ）
\end{itemize}
どの経路でも，最終段階で
\begin{itemize}
  \item 回転行列なら直交性（$R^\top R\approx I$）を確認
  \item クォータニオンなら正規化（$\|q\|=1$）
\end{itemize}
を行うと，数値誤差の蓄積を抑えられる。

\subsection*{2. 透視投影行列の計算}
画素ピッチ（1画素あたりの長さ）は
\[
p_x=\frac{33.6}{6720}=0.005\ \mathrm{mm/pixel},\quad
p_y=\frac{22.4}{4480}=0.005\ \mathrm{mm/pixel}.
\]
したがって焦点距離の画素単位は
\[
f_x=\frac{50}{0.005}=10000,\quad
f_y=\frac{50}{0.005}=10000.
\]
主点はセンサ中心より
\[
c_x=\frac{6720}{2}=3360,\quad c_y=\frac{4480}{2}=2240.
\]
スキュー $s=0$ とすると，内部パラメータ行列は
\[
K=
\begin{bmatrix}
10000 & 0 & 3360\\
0 & 10000 & 2240\\
0 & 0 & 1
\end{bmatrix}.
\]
カメラ座標系で外部パラメータを $[R|t]=[I|0]$ と置けば，透視投影行列は
\[
P=K[I|0]=
\begin{bmatrix}
10000 & 0 & 3360 & 0\\
0 & 10000 & 2240 & 0\\
0 & 0 & 1 & 0
\end{bmatrix}.
\]
よって3次元点 $\mathbf{X}=(X,Y,Z,1)^\top$ の画像座標は
\[
\lambda
\begin{bmatrix}
u\\v\\1
\end{bmatrix}
=
P\mathbf{X},\quad
u=\frac{10000X+3360Z}{Z},\ 
v=\frac{10000Y+2240Z}{Z}.
\]

\section*{考察}
回転表現はそれぞれ長所が異なる。回転行列は合成が容易だが冗長，オイラー角は直感的だがジンバルロックがある。Rodrigues（軸角）は最小3自由度で最適化に使いやすく，クォータニオンは補間（SLERP）と数値安定性に優れる。用途に応じて表現を切り替え，最終的に投影計算では行列形式へ統一するのが実装上有効である。

\end{document}
